% =====================================================================================
% Appendix C · Environments and reproduction.
%
% The detailed version of the front matter's short note. Kept separate rather than merged
% because the front matter must be readable in one sitting and this is a reference: an
% environment table, a notebook-to-chapter map, and the build pipeline.
%
% No number here is a result — the pymc version boundary and the dataset sizes are facts
% about the software and the files, not measurements — so no macro is cited.
% =====================================================================================
\chapter{Environments and reproduction}
\label{app:env}

Everything in this book is runnable, and this appendix is what it takes to run it. The short
version is in the front matter, under \emph{Reproducing this book}; what follows is the
reference.

% =====================================================================================
\section{Why there are two environments}
\label{app:env:split}

The single largest engineering decision behind the code is that it ships \emph{two} Python
environments rather than one, and the reason is not preference. As of writing, the PyMC ecosystem
is genuinely split across the \code{pymc 5} $\to$ \code{6} boundary: \code{pymc-bart},
\code{bambi} and \code{pathmc} require pymc~$\ge$~6, while \code{causalpy} and
\code{pymc-marketing} still pin pymc~$<$~6. No single environment satisfies both.

Three responses were available. Hold the whole book back to pymc~5, and lose the modern BART and
path-analysis work that two chapters are built on. Vendor patched forks, and hand the reader a
result they cannot reproduce from the public packages. Or ship two environments, and say so. The
third is what this book does, and it is stated here rather than buried in a lockfile because a
reader who tries to install one environment and run everything in it will otherwise conclude,
reasonably and wrongly, that something is broken.

\begin{center}
\footnotesize
\begin{tabular}{@{}l@{\hspace{1em}}c@{\hspace{1em}}l@{\hspace{1em}}p{0.34\linewidth}@{}}
\toprule
\textbf{Environment} & \textbf{pymc} & \textbf{Built by} & \textbf{Powers} \\
\midrule
\code{cmp-core} & $\ge$ 6 & \code{uv sync} &
  Notebooks 00--05, 07, 07b; the interactive apps; the package tests. \\[3pt]
\code{cmp-legacy} & $<$ 6 & \code{make env-legacy} &
  Notebooks 06, 08--11, 11b --- everything built on \code{causalpy} or
  \code{pymc-marketing}. \\
\bottomrule
\end{tabular}
\end{center}

The two live side by side in \code{.venv} and \code{.venv-legacy}, and neither is a subset of the
other. \code{cmp.estimators} imports \code{pymc} and \code{pymc-bart} \emph{lazily}, so the package
itself imports cleanly under both. Each notebook declares the kernel it needs, and the test suite
routes each notebook to the right interpreter, so \code{make test} drives both from a single
command. When \code{causalpy} and \code{pymc-marketing} move to pymc~6, the split collapses: delete
the legacy requirements file and bump the pins.

\begin{quote}\small
\code{uv sync} --- build the core environment.\\
\code{make env-legacy} --- build the legacy one.\\
\code{make kernels} --- register both as Jupyter kernels.
\end{quote}

% =====================================================================================
\section{Two environment variables}
\label{app:env:flags}

\subsection*{\code{CMP\_FAST} --- teaching mode}

Default \code{1}. Short chains, few trees, small placebo counts: every notebook runs in under about
two minutes, which is what makes a live rerun possible in front of a class and a full notebook test
suite possible in continuous integration.

\textbf{Every number in this book comes from \code{CMP\_FAST=0}} --- the full run: long chains, the
full seed loops, the full placebo distributions. The fast mode is a teaching device and nothing in
these pages was graded on it.

\subsection*{\code{CMP\_REAL} --- fetch the real data}

Default \code{0}. The book's default data is \emph{simulated}, for the reason \cref{app:dgp} exists:
a method has to be graded against a known truth before it can be trusted where there is none. It
also keeps the repository offline and deterministic.

Setting \code{CMP\_REAL=1} enables the cells that fetch a real public dataset over the network.
Nothing has to be supplied by the reader; the files are public, they are fetched rather than
vendored, and they are cached under \code{\textasciitilde/.cache/cmp} so that the download happens
once.

\begin{center}
\footnotesize
\begin{tabular}{@{}l@{\hspace{1em}}p{0.34\linewidth}@{\hspace{1em}}p{0.30\linewidth}@{}}
\toprule
\textbf{Chapter} & \textbf{Dataset} & \textbf{What replaces the planted truth} \\
\midrule
\cref{ch:foundations}, \cref{ch:ctrl} & LaLonde / NSW job training, in the
  Dehejia--Wahba composite &
  A published \emph{experimental} benchmark, cited (see \emph{A note on the numbers}). \\[3pt]
\cref{ch:geo} & Google's published geo experiment: 100 geos, one campaign window &
  A randomized, matched-pairs estimate --- unbiased by construction. \\[3pt]
\cref{ch:iv} & Criteo's advertising experiment (a one-time download of a few hundred
  megabytes) &
  A full-file Wald estimate, plus disjoint-fold stability. \\
\bottomrule
\end{tabular}
\end{center}

Chapter~3's real-data section runs on Hillstrom's randomized email test, which has no ground-truth
CATE --- real customers have one future each --- and is therefore graded on what a randomized
experiment \emph{can} grade: the average effect, and the realised profit of a targeting policy.

% =====================================================================================
\section{The notebooks, and the chapters they became}
\label{app:env:map}

Fourteen notebooks; thirteen chapters. The two real-data companions fold into their parent
chapters as an \emph{On real data} section rather than standing alone, because the point they make
is about the method, not about the dataset.

\begin{center}
\footnotesize
\begin{longtable}{@{}l@{\hspace{0.8em}}c@{\hspace{0.8em}}p{0.15\linewidth}@{\hspace{0.8em}}p{0.28\linewidth}@{}}
\toprule
\textbf{Notebook} & \textbf{Env} & \textbf{Chapter} & \textbf{The decision} \\
\midrule
\endhead
\code{00\_foundations} & core & \cref{ch:foundations}, \cref{ch:pvp} &
  Did the email work, and what is a posterior for? \\
\code{01\_uplift\_targeting} & core & Chapter~3 & Who should get the discount? \\
\code{02\_segment\_effects} & core & \cref{ch:seg} & Does the offer work better for some
  segments? \\
\code{03\_price\_elasticity} & core & \cref{ch:pe} & How should each region be priced? \\
\code{04\_funnel\_mediation} & core & \cref{ch:med} & Which funnel stage deserves the
  investment? \\
\code{05\_what\_to\_control\_for} & core & \cref{ch:ctrl} & Which variables belong in the
  model? \\
\code{06\_incrementality\_mmm} & legacy & Chapter~8 & Is the spend \emph{causing} the sales? \\
\code{07\_geo\_lift} & core & \cref{ch:geo} & Did the regional campaign work? \\
\code{07b\_geo\_lift\_real} & core & \cref{sec:sc:real} & The same question, on a real
  randomized geo experiment. \\
\code{08\_rollout\_did} & legacy & \cref{ch:did} & Did the loyalty rollout work? \\
\code{09\_threshold\_perk\_rdd} & legacy & \cref{ch:rdd} & Does the Gold perk retain
  people? \\
\code{10\_redesign\_its} & legacy & \cref{ch:its} & Did the site redesign help? \\
\code{11\_endogenous\_exposure\_iv} & legacy & \cref{ch:iv} & Ad exposure is self-selected ---
  what is its true effect? \\
\code{11b\_endogenous\_exposure\_real} & legacy & \cref{sec:iv:real} & The same question, on a
  real advertising experiment. \\
\bottomrule
\end{longtable}
\end{center}

\noindent\emph{The two \code{b} notebooks carry a \code{\_data} suffix on disk:}
\begin{quote}\footnotesize
\code{07b\_geo\_lift\_real\_data.ipynb}\\
\code{11b\_endogenous\_exposure\_real\_data.ipynb}
\end{quote}

% =====================================================================================
\section{Building the book}
\label{app:env:build}

\begin{quote}\small\ttfamily
make book
\end{quote}

which is \code{python book/build.py}, and which does four things in order:

\begin{enumerate}[leftmargin=*, itemsep=.3em]
  \item merges the notebooks' emitted result shards and writes one \code{\textbackslash newcommand}
        per number into the generated macro file;
  \item turns every emitted figure into a complete float --- caption, label and all --- so that no
        caption in this book was written anywhere but beside the computation it describes;
  \item runs the typesetter and the bibliography engine to produce the book; and
  \item produces the \emph{offprints}: one PDF per chapter, carrying the book's own chapter,
        equation and float numbers --- an offprint, not a second edition --- so that a reader
        handed \cref{ch:geo} alone finds its tables and equations numbered exactly as they are
        here.
\end{enumerate}

Two behaviours of that script are worth knowing because they look like defects and are not.

\textbf{It refuses to ship a PDF with a hole in it.} A chapter citing a number the notebooks did
not emit produces an undefined macro. The typesetter would happily continue and print a blank; the
build script reads the log, finds the undefined macro, and exits with an error naming it. A hole is
louder than a stale number, and the whole apparatus of \emph{A note on the numbers} rests on that
refusal.

\textbf{A fresh clone cannot build the book.} The generated directory holds nothing until the
notebooks have run, so the first step of any reproduction is to execute them at full quality
(\code{CMP\_FAST=0}), in both environments, with \code{CMP\_REAL=1} for the chapters that grade
themselves on real data. Only then does the book typeset. That the book will not compile against an
empty results directory is the same property, viewed from the other side, that guarantees no number
in it is stale.

\begin{quote}\small
\code{make test} --- executes every notebook headless, in the right environment, and runs the
package unit tests.\\
\code{make test-fast} --- the unit tests alone.
\end{quote}

\begin{remark}[There is no index]
This book has no index, and the omission is deliberate rather than pending. An index worth having
requires index entries placed in the prose --- \emph{this} occurrence of ``parallel trends'' is the
one a reader should be sent to, not the eleven others --- and entries have to be authored, chapter
by chapter, by someone who knows which occurrence is the definition. An index generated
mechanically from a keyword list sends the reader to the wrong page with the full authority of
typography, and is worse than no index at all. The table of contents is detailed to the section, the
list of figures and the list of tables are complete, and the cross-references are live. When the
entries are written, the pipeline gains one pass and the index appears; until then, this is the
honest state.
\end{remark}
